\uuid{lSza}
\exo7id{1058}
\titre{exo7 1058}
\auteur{legall}
\organisation{exo7}
\datecreate{1998-09-01}
\isIndication{false}
\isCorrection{false}
\chapitre{Matrice}
\sousChapitre{Propriétés élémentaires, généralités}
\cpge{mpsi-mp2i.calcul-matriciel-systemes.e, pcsi.determinants.c, ptsi.determinants.a, mpsi-mp2i.espaces-vectoriels.b, pcsi.espaces-vectoriels.b, ptsi.espaces-vectoriels.a, bcpst1.matrices, pc.complements-espaces-vectoriels.b, psi.complements-espaces-vectoriels.b, pt.complements-algebre-lineaire.b}
\module{Algèbre}
\niveau{L1}
\difficulte{}

\contenu{
\texte{
Soient $  A  $ et $  B \in M_n(\mathbb{K} )  $ deux matrices triangulaires
sup\'erieures.
}
\begin{enumerate}
    \item \question{Montrer (en calculant les coefficients) que $  AB  $ est triangulaire sup\'erieure.}
    \item \question{Soit $  \varphi   $ un endomorphisme bijectif de $  \mathbb{K} ^n  $ et $  F  $
un sous-espace vectoriel de $  \mathbb{K} ^n  $ tel que $  \varphi (F) \subset F  .$ Montrer que
 que $  \varphi ^{-1}(F) \subset F  .$}
    \item \question{En d\'eduire une nouvelle d\'emonstration de 1. Montrer que si $  A  $ est inversible, $  A^{-1}  $ est triangulaire sup\'erieure.}
\end{enumerate}
}
